% SPDX-License-Identifier: Apache-2.0
\section{Units, Processes and Parallelism}
\label{sec:e-units}

A unit is a struct that implements \code{crate::comp::Unit}. The trait
has one method.

\begin{lstlisting}[caption={The unit trait. Inputs and outputs are type parameters, so one struct may implement it more than once.},label={lst:module}]
#[allow(async_fn_in_trait)]
pub trait Unit<In, Out> {
    async fn run(&mut self, inputs: In, outputs: Out);
}
\end{lstlisting}

Inputs and outputs are type parameters rather than associated types, and
that choice buys something concrete: one struct may implement the trait
more than once, with different shapes, so a unit that works on 32-bit and
64-bit streams is one implementation and not two.

\subsection{Parallelism is stated, not implied}

Synthesis calls \code{run}. A unit with more than one concurrent process
starts one async function per process and joins them.

\begin{lstlisting}[caption={Parallelism written down. A unit with one process writes no join.},label={lst:join}]
impl Unit<Ins, Outs> for DualPort {
    async fn run(&mut self, _i: Ins, _o: Outs) {
        join2(self.port_a(), self.port_b()).await;
    }
}
\end{lstlisting}

This replaces a rule with a construct. TxHDL said that several sequence
blocks in one module run concurrently, which was a property of having
declared them. It also contradicted itself: the specification stated that a
module holds one sequence, while three of its own modules declared
two~\cite{unification}. Here there is no rule to contradict. Parallelism is
a \code{join}, written where it happens, and a unit with one process does
not write one.

\subsection{Processes cannot take a mutable borrow}

The obvious way to write the two processes does not compile, and finding
that out cost ten lines.

\begin{lstlisting}[caption={What does not compile, and the error it gives. This is the code a Rust programmer writes first.},label={lst:e0499}]
async fn port_a(&mut self) { /* ... */ }
async fn port_b(&mut self) { /* ... */ }

// error[E0499]: cannot borrow `*self` as mutable
//               more than once at a time
join2(self.port_a(), self.port_b()).await;
\end{lstlisting}

Hardware processes share state: they drive registers. Rust permits one
mutable borrow at a time. The two facts collide directly.

The repair is that a register is a cell, and a process takes a shared
reference.

\begin{lstlisting}[caption={The repair. A register is a cell several processes reach, which is what a register is. Reading and driving it are plain; the wait is the edge.},label={lst:reg}]
pub struct Reg<T: Copy, C: Clock = DefaultClock> { .. }

impl<T: Copy, C: Clock> Reg<T, C> {
    pub fn get(&self) -> T;      // what the last edge latched
    pub fn set(&self, x: T);     // taken at the end of the step
}

async fn port_a(&self) {
    loop {
        DefaultClock::rising().await;    // the wait
        self.hits_a.set(self.hits_a + 1);  // a read at that edge
    }
}
\end{lstlisting}

\code{run} still takes \code{\&mut self}, and the processes it starts take
\code{\&self}. That compiles, and it also describes the hardware better: a
register is a thing several processes reach, which is what a register is.

\subsection{A process waits, then reads, and it loops}
\label{sec:e-regread}

Listing~\ref{lst:reg} makes three claims that earlier drafts did not,
and the first took three drafts of its own to get right.

\textbf{Waiting and reading are different things.} A process waits for
an event: the next edge of a clock, \code{C::rising()} or
\code{C::falling()}; a transaction on a channel, \code{rx.wait()}; or
an edge at which a condition on state holds, \code{until(C::rising,
|| ..)} and the same with \code{falling}. Once its event has happened it reads
state, plainly. \code{Reg::get} returns what the last edge latched and
waits for nothing, and so does a wire's \code{In::get}. The two drafts
before this one made the register read the wait, \code{get} being
\code{async}, on the argument that the \code{.await} on a read was the
mark the lowering turns into a register. It was also the only wait
there was, so a process that wanted to react to a channel and then
look at a pointer had no way to say it: a second \code{.await} for the
pointer was a second edge, and a group of the two gave the pointer
from before the channel's edge. Separating the wait from the read
removed the problem and left the mark where it belongs: the edge is
the cycle boundary, and everything read after it is what the hardware
reads at that edge.

\textbf{Driving is deferred.} \code{set} states the next value, and the
register takes it at the end of the step, so a read after a drive in
the same iteration still sees what the edge latched. That is what a
register does, and a testbench that reads a register between steps
sees the committed value without a second kind of read.

\textbf{A process loops.} A unit runs for as long as the clock does, so
every leaf \code{run} is a \code{loop}, and each iteration contains a
wait. A parent whose \code{run} only joins its children needs no loop
of its own: the children loop, the join never completes, and the
parent is continuous through them.

The prototype enforces this in the only way a simulation can. Every
\code{.await} is a cycle boundary: a process that has crossed an edge
at this step cannot cross it again, so a second wait written after a
first waits for the next edge. Waits that share an edge say so:

\begin{lstlisting}[caption={Two events at one edge are written inside \code{parallel!}. Written one after the other, they would be two edges, as any two awaits are.},label={lst:parallel}]
let (x, y) = parallel!(rx_a.wait(), rx_b.wait()).await;
\end{lstlisting}

\code{parallel!} polls its futures together and returns their outputs
as a tuple, and it is the one thing the executor treats specially: once
one wait of the group has crossed the edge at this step, the others in
the group are free, each once. \code{join2} and \code{parallel!} are the
two sides of one idea: the first gives each of its futures a process,
the second keeps them in one, and a unit uses both. Processes are told
apart by their waker, which \code{join2} and \code{join\_all} give each
child fresh. A loop that awaits nothing never yields and spins. That is
the right behaviour, because such a loop is a combinational loop in
the hardware, which is the thing that was wrong.

A key timing subtlety arises from loop execution order. State sampled
before an \code{await} statement at the top of a loop reflects the
preceding clock step, prior to the commitment of that step's driven values.
If an arbiter samples its round-robin priority before waiting for request
assertion, it samples stale cycle state, causing double grants. Priority
evaluation must occur strictly after the wait resumes. For this reason,
arbitration blocks on an \code{until} condition combining input requests
via \code{||}, sampling winner priority immediately upon waking. The
\code{select!} macro provided by the library does not block; it evaluates
combinational pattern matches, as described in Section~\ref{sec:e-select}.

\begin{pitfall}
This is the constraint most likely to surprise somebody writing a unit for
the first time, because the code that fails is the code a Rust programmer
would write first. It should be in the library's documentation and in the
error message, not only in a paper.
\end{pitfall}


\subsection{A unit that contains units}
\label{sec:e-hierarchy}

A submodule is a field. A bus that joins two submodules is a field too,
owned by whichever unit contains both of its ends.

\begin{lstlisting}[caption={Submodules and the wire between them are fields of the parent. Each child is generic over its own configuration, not the parent's.},label={lst:hier}]
pub struct Top<TC: TopConfig> {
    pub producer: Producer<TC::P>,
    pub consumer: Consumer<TC::C>,
    pub link: DataBus,
    pub pes: [Pe; 4],
}
\end{lstlisting}

\subsection{Configurations nest}
\label{sec:e-nestedconfig}

Section~\ref{sec:e-binding} gives a design one configuration and makes
every unit generic over it. That holds while every choice is global. It
stops holding as soon as a submodule has choices of its own, and
Listing~\ref{lst:hier} is where that shows.

A configuration is a tree. A unit declares the configuration \emph{it}
needs, and a parent's configuration names its children's rather than
holding their fields.

\begin{lstlisting}[caption={Each unit declares what it needs. The parent names its children's configurations and never mentions their knobs.},label={lst:nestedcfg}]
pub trait ProducerConfig {
    type Gen: Generator;
    const DEPTH: usize;
}

pub trait ConsumerConfig {
    const DEPTH: usize;        // the same name, a different unit
    const CHECKSUM: bool;
}

pub trait TopConfig {
    type P: ProducerConfig;
    type C: ConsumerConfig;
    const NAME: &'static str;
}
\end{lstlisting}

Two things follow, and both were the reason to do it.

\textbf{Two children may want a knob of the same name.} Both
\code{DEPTH}s above exist, differ, and never meet, because each belongs
to its own trait. A flat configuration would have had to rename one of
them after the fact, which is a rename in a design that had nothing wrong
with it.

\textbf{A child's configuration is reusable.} It is written once and
named by whichever parents want it, so two builds may differ in one
subsystem and share the rest.

\begin{lstlisting}[caption={Two builds. The consumer's configuration is written once and used by both; only the producer differs.},label={lst:nestedbuilds}]
pub struct CheckedConsumer;
impl ConsumerConfig for CheckedConsumer {
    const DEPTH: usize = 64;
    const CHECKSUM: bool = true;
}

pub struct Fpga;
impl TopConfig for Fpga {
    type P = FastProducer;
    type C = CheckedConsumer;
    const NAME: &'static str = "fpga";
}

pub struct Tiny;
impl TopConfig for Tiny {
    type P = SmallProducer;
    type C = CheckedConsumer;   // reused, unchanged
    const NAME: &'static str = "tiny";
}
\end{lstlisting}

Nesting costs nothing at elaboration. A child's constant is still a
compile-time value reached through the parent, so it may still size an
array:

\begin{lstlisting}[caption={A child's constant, reached through the parent, still sizes an array. This is the long spelling.},label={lst:nestedconst}]
pub const DEPTH: usize =
    <<Fpga as TopConfig>::P as ProducerConfig>::DEPTH;
pub static BUF: [u32; DEPTH] = [0; DEPTH];
\end{lstlisting}

The spelling in Listing~\ref{lst:nestedconst} is the one real cost of
nesting: reaching two levels down needs a fully qualified path, and three
levels reads worse again. The repair is one type alias per level of the
tree, written once beside the trait, so it serves every build rather than
one.

\begin{lstlisting}[caption={The repair. A generic alias per level, and a per-build alias where one build is being talked about. Both compile in probe~17.},label={lst:nestedalias}]
pub type ProducerOf<T> = <T as TopConfig>::P;
pub type ConsumerOf<T> = <T as TopConfig>::C;

// One qualification, not two, for any build.
pub const DEPTH: usize =
    <ProducerOf<Fpga> as ProducerConfig>::DEPTH;

// None at all, where one build is meant.
pub type FpgaProducer = ProducerOf<Fpga>;
pub const D: usize = <FpgaProducer as ProducerConfig>::DEPTH;
pub static BUF: [u32; D] = [0; D];
\end{lstlisting}

A derive on the configuration trait could emit the per-level aliases,
which would make the long spelling something nobody writes.

Wiring is not a stored back-reference. A child is handed the ports it
needs when its \code{run} is called, and that is what the inputs and
outputs of \code{Unit} are for.

\begin{lstlisting}[caption={Instantiating two submodules and wiring them. Both children receive the same bus; neither holds a reference to the other.},label={lst:hier-run}]
impl<TC: TopConfig> Unit<(), ()> for Top<TC> {
    async fn run(&mut self, _i: (), _o: ()) {
        join2(
            self.producer.run((), ProducerOut { out: &self.link }),
            self.consumer.run(ConsumerIn { inp: &self.link }, ()),
        ).await;
    }
}
\end{lstlisting}

\subsection{Why the borrow rule does not bite here}

Section~\ref{sec:e-units} showed that two processes belonging to the
same unit cannot concurrently hold \code{\&mut self}. Listing~\ref{lst:hier-run}
appears to permit this and compiles cleanly.

The mechanisms are fundamentally distinct. A process method borrows the
entire struct reference, causing conflicting borrows across processes. In
contrast, a submodule is an individual struct \emph{field}. The Rust borrow
checker permits simultaneous mutable borrows of disjoint fields, allowing
\code{\&mut self.producer} and \code{\&mut self.consumer} to operate
concurrently. The interconnect bus between them is borrowed immutably by
both endpoints, accurately reflecting shared wiring.

\begin{rulebox}[Where interior mutability is needed]
State that two processes of one unit both drive. Not submodules, which are
disjoint fields, and not the buses between them, which are shared by
design.
\end{rulebox}

\subsection{Arrays of instances}

An array of identical submodules is an array, and joining them is a join
over an iterator.

\begin{lstlisting}[caption={An instance array. The replication count is a Rust array length, so a config can set it.},label={lst:hier-array}]
join_all(
    self.pes.iter_mut()
        .enumerate()
        .map(|(i, pe)| pe.run(i as u32, ()))
).await;
\end{lstlisting}

TxHDL wrote this as \code{instance pe[16]} with a \code{connect} loop
beside it. Here the replication is an array length and the connection is
what each element is passed, so a config may set the count through an
associated constant and nothing else changes.
